\section{Hamilton 模型与固定核例}
\label{sec:rank-deficient-hamilton}

\subsection{秩亏 Hamilton 系统的完整近端收敛}

考虑秩亏设计矩阵 $(1,0)$ 对应的幂次回归损失
\[
 f(x)=\frac23|x_1|^{3/2},\qquad x=(x_1,x_2)\in\mathbb R^2,
\]
以及动能 $k(p)=\frac23\|p\|^{3/2}$。取阻尼 $\Gamma=99$，
全隐式 Hamilton 更新为
\begin{equation}
\label{eq:hd-implicit}
 x^{j+1}-x^j=\nabla k(p^{j+1}),\qquad
 p^{j+1}-p^j=-\nabla f(x^{j+1})-99p^{j+1}.
\end{equation}
记
\[
 g(x)=\bigl(\operatorname{sgn}(x_1)\sqrt{|x_1|},0\bigr),
 \qquad
 h(p)=\begin{cases}p/\sqrt{\|p\|},&p\ne0,\\0,&p=0.\end{cases}
\]
在 $X=\mathbb R^2\times\mathbb R^2$ 上定义
\[
 F(x,p)=(-h(p),g(x)+99p).
\]
于是 \eqref{eq:hd-implicit} 等价于 $z^{j+1}=J_F(z^j)$，
其中 $z^j=(x^j,p^j)$。其零点集为
\[
 S=F^{-1}(0)=\{((0,s),0):s\in\mathbb R\}.
\]

\begin{proposition}
\label{prop:rank-deficient-hamilton-rleb}
令 $U=X$、$R=10^{-6}$、$U_R=\{z:d(z,S)\leq R\}$，并定义
完整近端输入图块
\[
 \mathcal G_R=\{(y,F(y)):y+F(y)\in U_R\}.
\]
则 $J_F$ 在整个 $X$ 上单值且处处有定义，
$T_{\mathcal G_R}=J_F$ 在 $U_R$ 上成立。
该图块满足引理~\ref{lem:RL-J} 的全对近端能量条件，
并通过最近零点锚验证假设~\ref{ass:rleb-ppa}，可取
\[
 \lambda=1,\qquad \gamma=\frac12,\qquad L=\frac12,
 \qquad \bar t=1,\qquad \psi(t)=3t^2,
 \qquad \kappa=\frac14.
\]
这里的 RL 条件对图块中输入差不超过 $R$ 的全部点对成立，
相应的允许关系取完整近端 $\mathcal J=J_F$。
\end{proposition}

\begin{proof}
将近端输入记为 $(a,b)$、输出记为 $(x,p)$，则近端方程为
\begin{equation}
\label{eq:hd-proximal-equations}
 a=x-h(p),\qquad b=g(x)+100p.
\end{equation}
由于 $k^*(v)=\|v\|^3/3$，函数
\[
 u\longmapsto k^*(u-a)+\frac1{100}f(u)
                         -\frac1{100}\langle b,u\rangle
\]
严格凸且强制，因而有唯一极小点 $x$。
其一阶条件与 $p=\nabla k^*(x-a)=\|x-a\|(x-a)$
恰好给出 \eqref{eq:hd-proximal-equations}。
故 $J_F$ 在 $X$ 上处处存在且唯一。
对每个 $U_R$ 中的输入，唯一输出及其图像点都被 $\mathcal G_R$ 保留；
此外，每个 $s\in S$ 满足 $(s,0)\in\mathcal G_R$。
最近零点具体为 $((0,a_2),0)$，因此输入覆盖和图块最近零锚均成立。

先证明：当 $(a,b)\in U_R$ 时，
\begin{equation}
\label{eq:hd-momentum-small}
 \|p\|\leq R.
\end{equation}
记 $P=\|p\|$。输入条件给出 $|a_1|\leq R$、$\|b\|\leq R$。
若 $p_1g_1(x_1)\geq0$，则展开 $\|g(x)+100p\|^2$ 可得
$100P\leq\|b\|$。若 $p_1g_1(x_1)<0$，则
$p_1$ 与 $x_1=a_1+p_1/\sqrt P$ 异号，故
\[
 |p_1|\leq R\sqrt P,\qquad |p_2|=|b_2|/100\leq R/100,
 \qquad P^2\leq R^2P+R^2/10000.
\]
若 $P>R$，则 $P^2-R^2P>R^2(1-R)>R^2/10000$，矛盾。
这证明了 \eqref{eq:hd-momentum-small}，并包括 $p=0$ 的情形。

其次，取任意两个 $U_R$ 中的近端输入。
令 $\beta=\sqrt2/4$，则附录中的
引理~\ref{lem:hamilton-momentum-difference} 给出原方程的有限差分界
\[
 |p_1-\widetilde p_1|
 \leq\max\left\{
       \frac{|b_1-\widetilde b_1|}{100},\,
       2\sqrt R|a_1-\widetilde a_1|
       +\frac{\beta}{100}|b_2-\widetilde b_2|
      \right\}.
\]
结合 $p_2=b_2/100$，得到
\begin{equation}
\label{eq:hd-momentum-lipschitz}
 \|p-\widetilde p\|
 \leq\left(4R+\frac{2+\beta^2}{10000}\right)^{1/2}
        \|(a,b)-(\widetilde a,\widetilde b)\|
 \leq0.02\|(a,b)-(\widetilde a,\widetilde b)\|
\end{equation}
这一论证只使用 $g_1$
的严格递增性；在 $p=0$、$x_1=0$ 及跨越这些点时仍成立，
无需平滑正则化。

现取两输入之差的范数 $\delta\le R$，并记
\[
 A_0=a-\widetilde a,\quad B_0=b-\widetilde b,\quad
 P_0=p-\widetilde p,\quad H_0=h(p)-h(\widetilde p).
\]
由引理~\ref{lem:hamilton-h-holder} 及式~\eqref{eq:hd-momentum-lipschitz}，
\[
 \|P_0\|\le0.02\delta,\qquad
 \|H_0\|^2\le2\|P_0\|\le0.04\delta.
\]
令 $y=(x,p)$、$\widetilde y=(\widetilde x,\widetilde p)$。
由原近端方程，输出差与原图值差分别为
$(A_0+H_0,P_0)$ 和 $(-H_0,B_0-P_0)$。
因此直接展开得到平方能量
\[
 \begin{aligned}
 E&:=\|y-\widetilde y\|^2+\|F(y)-F(\widetilde y)\|^2\\
 &=\delta^2+2\|H_0\|^2+2\|P_0\|^2
       +2\langle A_0,H_0\rangle-2\langle B_0,P_0\rangle\\
 &\le0.08\delta+0.4\delta^{3/2}+1.0408\delta^2\\
 &\le0.0804010408\delta
 \le\frac12(\delta^2+L^2\delta).
 \end{aligned}
\]
最后一步使用了 $R=10^{-6}$ 和 $L=1/2$；$\delta=0$ 时两图点相同。
引理~\ref{lem:RL-J} 遂给出 $\mathcal G_R$ 上同一尺度的全对 RL 条件。
每个实际输出都可与输入最近零点组成图块内点对，
其输入差正是输入到 $S$ 的距离，故条件 (P2) 成立。

再令 $(x,p)=J_F(a,b)$、$(a,b)\in U_R$，
使用原算子 $F$ 的真残差
\[
 \rho=r_F(x,p)=\|F(x,p)\|,\qquad
 \rho^2=\|p\|+\|g(x)+99p\|^2.
\]
从 \eqref{eq:hd-momentum-small} 得到
\[
 \|p\|\leq\rho^2,\qquad
 \|g(x)\|\leq\rho+99\|p\|
             \leq(1+99\sqrt R)\rho.
\]
于是
\[
 \begin{split}
 d((x,p),S)&=\sqrt{|x_1|^2+\|p\|^2}\\
 &\leq |x_1|+\|p\|
  =\|g(x)\|^2+\|p\|\\
 &\leq\bigl((1+99\sqrt R)^2+1\bigr)\rho^2
  <3\rho^2\qquad(\rho>0).
 \end{split}
\]
当 $\rho=0$ 时输出属于 $S$，同一非严格误差界仍成立。
实际上 $g(x)+99p=b-p$，故
$\rho^2\leq R+4R^2<1$，整个输出块都在所取残差区间内。
再由 $(R+L\sqrt R)/2=0.0002505<1=\bar t$，
验证了 (P3)。

最后，对任意 $0<r\leq R$，有
\[
 \frac{\psi((r+L\sqrt r)/2)}r
 =\frac34\left(\sqrt r+\frac12\right)^2
 \leq\frac34(0.001+0.5)^2
 =0.18825075<\frac14.
\]
这证明 (P4) 的收缩条件。
此外，$B(r)=(r+L\sqrt r)/2$ 在几何距离尺度上的和有限，
故 (P4) 的可和性也成立。
\end{proof}

这里 $U=X$，故 $d(z^0,X\setminus U)=+\infty$，
定理~\ref{thm:rleb-ppa-general} 的留域条件对每个 $z^0\in U_R$ 自动成立。
因此，对每个满足 $d_0:=d(z^0,S)\leq10^{-6}$ 的初值，
迭代 \eqref{eq:hd-implicit} 在 $U_R$ 内无限延拓，并收敛到
某个 $z^\infty\in S$；其距离与尾项具有上界
\[
 d(z^j,S)\leq4^{-j}d_0,\qquad
 \|z^j-z^\infty\|
 \leq\frac23\,4^{-j}d_0+\frac12\,2^{-j}\sqrt{d_0}.
\]
这些是定理给出的统一上界，不主张是每条轨道的精确渐近速度。
结论允许非零初始动量，并控制整条位置--动量轨道的极限点。

本例在完整局部图块上的最大 RL 指数为 $\gamma=1/2$。
秩亏方向给出一族显式输入：令
$w_t=((0,0),(0,t))$、$w_0=0$，其中 $0<t\le R$，则
\[
 J_F(w_t)=\bigl((0,\sqrt{t/100}),(0,t/100)\bigr),\qquad
 \|J_F(w_t)-J_F(w_0)\|\ge\frac1{10}\sqrt t.
\]
因此完整预解式在零点附近不局部 Lipschitz。
由于两输入的距离为 $t$，同一点对的平方能量也排除任何
$\gamma>1/2$ 的 RL 条件；其详细计算见
命题~\ref{prop:hamilton-sharp-exponents}。
其中还证明了真残差误差界的最大幂次为 $2$，
并给出一条可显式求解、快于统一上界的轨道。

另一个可直接检验的区别是能量配对条件。
若要求以统一常数控制
$|\langle\nabla f(x),\nabla k(p)\rangle|$
与 $f(x)+k(p)$ 的比值，则本例不满足该条件：取
$x=p=(t,0)$、$t\downarrow0$，有
\[
 \frac{|\langle\nabla f(x),\nabla k(p)\rangle|}{f(x)+k(p)}
 =\frac{t}{\frac43t^{3/2}}
 =\frac3{4\sqrt t}\longrightarrow\infty.
\]
因此，要求这一配对条件的收敛命题不能直接用于本例。
